
LLM trustworthiness frameworks --- TrustLLM, HELM, MultiTrust --- treat safety, fairness, truthfulness, robustness, privacy, and machine ethics as independent axes requiring separate evaluation. This paper examines whether that independence assumption is empirically warranted. Computing partial Spearman rank correlations across TrustLLM's 16-model $\times$ 6-dimension benchmark data and controlling for model scale (log parameter count) and RLHF alignment status, we find that 8 of 15 dimension pairs are statistically significant at a Bonferroni-corrected threshold ($\alpha$ = 0.0033), with adversarial robustness as the sole RLHF-insensitive dimension while the remaining five dimensions form a correlated cluster ($\rho$ up to 0.971) shaped by RLHF preference training. A within-family LLaMA-2 natural experiment provides directional mechanistic evidence: RLHF fine-tuning jointly increases safety (mean $\Delta$ = +0.639) and machine ethics (mean $\Delta$ = +0.424) at 7B, 13B, and 70B scale. A minimum spanning tree of the correlation geometry identifies {truthfulness, fairness, privacy} as a sufficient 3-dimension evaluation set with mean bootstrap edge frequency 0.917 (1000 resamples) --- a 50\% reduction from the standard 6-dimension protocol. The predicted cluster boundary (RLHF-sensitive {safety, ethics} vs. RLHF-insensitive {robustness, privacy}) is partially refuted: privacy is the strongest co-mover with safety ($\rho$ = 0.971), indicating it is RLHF-sensitive against prior expectation. Adversarial robustness alone fails to co-move with the RLHF-shaped cluster, though the full partial correlation ($\rho = -0.188$, p = 0.519) does not reach significance at n = 16. These findings suggest trustworthiness evaluation be treated as a geometry problem rather than a dimension-counting problem.

